\documentclass[11pt,a4paper]{article}

\usepackage{geometry}
\geometry{margin=1in}

\usepackage{amsmath,amssymb}
\usepackage{graphicx}
\usepackage{booktabs}
\usepackage{hyperref}
\usepackage{enumitem}
\usepackage{array}
\usepackage{float}

\title{\textbf{Progress Report: System-Level Evaluation of Rust and Python for Machine Learning}}
\author{Project Elective}
\date{\today}

\begin{document}

\maketitle

\section{Overview of the Project}

This project studies the use of \textbf{Rust} as an alternative systems language for machine learning workflows traditionally implemented in \textbf{Python}.  
Rather than focusing on state-of-the-art model performance, the emphasis is on:

\begin{itemize}
    \item feasibility of end-to-end ML workflows,
    \item system stability and reproducibility,
    \item developer experience and DevOps complexity,
    \item deployment and operational characteristics.
\end{itemize}

To ensure clarity and rigor, the work is organized into \textbf{two clearly separated experimental tracks}.

---

\section{Project Structure: Two-Track Evaluation}

The project consists of the following two tracks:

\subsection*{Track 1: Training-Based Systems Evaluation}
This track compares \textbf{machine learning training pipelines} implemented in:
\begin{itemize}
    \item PyTorch (Python), and
    \item Burn (Rust).
\end{itemize}

The goal is to evaluate training feasibility, stability, compile-time guarantees, and DevOps impact, rather than raw training speed.

\subsection*{Track 2: Inference-Based DevOps Evaluation}
This track compares \textbf{production-style inference services} implemented in:
\begin{itemize}
    \item Python-based ONNX inference, and
    \item Rust-based ONNX inference.
\end{itemize}

The focus is on deployment, security, containerization, CI/CD behavior, and runtime efficiency.

Each track is designed to answer a distinct research question while remaining complementary.

---

\section{Machine Learning Tasks Considered}

To ensure coverage of diverse ML workloads, the following tasks are identified:

\begin{itemize}
    \item \textbf{Text Classification}: Dataset to be finalized.
    \item \textbf{Image Classification}: MNIST dataset.
    \item \textbf{Credit Score Assignment}: Supervised classification task.
    \item \textbf{Multi-Objective Machine Learning}: Brain Tumor dataset with a MOML formulation.
    \item \textbf{Fine-Tuning Task}: BERT-based classification (ANLP Assignment 1), with optional LoRA / QLoRA.
    \item \textbf{Autoregressive Decoding}: Experiments using the Burn framework.
\end{itemize}

At the current stage, the \textbf{MNIST image classification task has been fully implemented}.  
The corresponding training code is available in the project GitHub repository.

---

\section{Related Work}

The following research papers are being used to guide experimental design and evaluation:

\begin{itemize}
    \item \url{https://ieeexplore.ieee.org/document/11126113}
    \item \url{https://ieeexplore.ieee.org/document/11261485}
    \item \url{https://ieeexplore.ieee.org/document/11212348}
    \item \url{https://www.ijsred.com/volume8/issue2/IJSRED-V8I2P143.pdf}
\end{itemize}

---

\section{Code Repository and Current Status}

Project repository:
\begin{center}
\url{https://github.com/Abhinav-Kumar012/Rust_Python_ML_PE.git}
\end{center}

Current progress includes:
\begin{itemize}
    \item MNIST training pipeline implemented
    \item PyTorch baseline established
    \item Initial Rust (Burn) training setup completed
\end{itemize}

---

\section{Track 1: Training-Based Systems Evaluation}

\subsection{Objective}

The objective of this track is to answer the following research question:

\begin{quote}
\textit{Can Rust realistically support end-to-end machine learning training pipelines, and what system-level trade-offs does this introduce compared to PyTorch?}
\end{quote}

This track explicitly avoids speed-centric benchmarking and instead focuses on system behavior.

---

\subsection{Frameworks Compared}

\subsubsection{PyTorch (Baseline)}
\begin{itemize}
    \item Language: Python
    \item Training maturity: Very high
    \item Ecosystem: Extensive
\end{itemize}

\subsubsection{Rust (Burn)}
\begin{itemize}
    \item Language: Rust
    \item Training maturity: Emerging
    \item Design: Idiomatic Rust, native training support
\end{itemize}

---

\subsection{Experimental Controls}

\textbf{Fixed Across Both Implementations}
\begin{itemize}
    \item Dataset splits
    \item Number of epochs
    \item Batch size
    \item Optimizer type
    \item Learning rate
    \item Hardware
\end{itemize}

\textbf{Allowed Differences}
\begin{itemize}
    \item Internal kernel implementations
    \item Graph execution model
    \item Memory management
\end{itemize}

---

\subsection{Metrics Collected}

\begin{itemize}
    \item Training time per epoch (reported cautiously)
    \item Loss curves and convergence behavior
    \item Runtime failures and numerical stability
    \item Reproducibility across runs
    \item Environment setup and build complexity
    \item Dependency footprint and artifact size
\end{itemize}

---

\section{Track 2: Inference-Based DevOps Evaluation}

\subsection{Objective}

The objective of this track is to compare \textbf{deployment, security, and operational characteristics} of Python-based and Rust-based ML inference services executing the same ONNX model.

---

\subsection{Inference Services Compared}

\textbf{Python Service}
\begin{itemize}
    \item FastAPI + Uvicorn
    \item ONNX Runtime (Python)
\end{itemize}

\textbf{Rust Service}
\begin{itemize}
    \item Axum / Actix
    \item burn-rs
\end{itemize}

Both services expose identical inference endpoints and return identical outputs.

---

\subsection{Evaluation Dimensions}

\begin{itemize}
    \item CI/CD build behavior
    \item Container image size and layering
    \item Cold-start latency
    \item Inference latency and throughput
    \item Resource utilization
    \item Security and supply-chain surface
\end{itemize}

---

\section{Upcoming Work}

The following tasks are planned for the next phase of the project:
s
\begin{itemize}
    \item Develop production-style inference services for both Python and Rust.
    \item Write Dockerfiles for Python and Rust inference services.
    \item Set up Jenkins-based CI pipelines for inference, including build, test, containerization, and security scanning.
\end{itemize}


\end{document}
